% ============================================================================
% Article paragraph: the SA variant used in the paper (Taillard-style
% plateaus, instance-aware absolute temperatures). No grid-search details
% here — calibration values belong to the experiments section.
% Verified against metacarp/recocido_simulado.py (T_init = 20·d_max/n,
% T_end = 20·d_max/n², L = n²). Generated: 2026-07-15.
% ============================================================================
In this paper we applied the SA algorithm with the following modifications,
following the design guidelines of \cite{taillardDesignHeuristicAlgorithms2023}.
First, unlike the textbook SA that decreases the temperature after every
single iteration, this variant holds the temperature constant over a plateau
of iterations --- a Markov chain --- before reducing it, as suggested in
\cite{taillardDesignHeuristicAlgorithms2023}; the plateau length is adapted
to the instance as $L = n^2$, where $n$ is the number of required edges
(tasks), so that larger instances are granted proportionally longer chains
at each temperature. Second, the algorithm requires absolute starting and
ending temperatures, which we also derive from the instance:
$T_{\mathrm{init}} = 20\, d_{\mathrm{max}}/n$ and
$T_{\mathrm{end}} = 20\, d_{\mathrm{max}}/n^{2}$, where $d_{\mathrm{max}}$
is the largest finite entry of the all-pairs shortest-path (deadheading)
distance matrix. Since the cost change $\Delta$ of a random move is of the
order of a few deadheading distances, scaling $T_{\mathrm{init}}$ by
$d_{\mathrm{max}}/n$ makes the initial acceptance probability of worsening
moves comparable across instances of very different size and cost scale.
The temperature decreases geometrically between plateaus,
$T \leftarrow \alpha T$, and we maintain the usual meaning of the cooling
parameter $\alpha$; note that with these choices
$T_{\mathrm{end}}/T_{\mathrm{init}} = 1/n$, so the number of temperature
levels, $\lceil \ln n / \ln(1/\alpha) \rceil$, grows logarithmically with
the instance size --- the schedule lengthens automatically on larger
instances without introducing new parameters. Moves are drawn from the nine
neighborhood operators shared by this study
(Table~\ref{tab:neighborhood-operators}), selecting the inter-route group
with probability $\max(p_{inter}, 0.8)$ under capacity violation and
$p_{inter}$ otherwise, and candidates are compared through the penalized
objective $f(s) = c(s) + \omega \cdot \mathrm{viol}(s)$, where
$\mathrm{viol}(s)$ is the total capacity excess. Finally, the variant
incorporates a reheating mechanism that restores the temperature to
$\rho\, T_{\mathrm{init}}$ after a patience of non-improving temperature
levels, described in detail below; the search terminates when
$T \leq T_{\mathrm{end}}$, when consecutive reheats no longer improve the
incumbent, or when the evaluation/time budget is exhausted. A pseudocode is
shown in Algorithm~\ref{alg:sa_simple}.
